\section{The factorial-moment proof}\label{sec:factorial}
We now prove Theorem~\ref{fac:main}. The cycle estimate controls sparse injections and very large representation dimensions. The remaining dimensions require a second density estimate, obtained after replacing fixed-size middle blocks by uniform permutations. We prove both density assertions before making the representation split.

For the recursive description number the coordinates so that the palindrome uses layers $d,\ldots,1,1,\ldots,d$. A node of size $2^t$ varies in its low $t$ coordinates. This reversal of coordinate names does not change its law up to conjugation. Thus $K_N=T_N^*T_N$, consistently with Section~\ref{sec:introduction}, and $\|K_N(\lambda)\|=\|T_N(\lambda)\|^2$.

\subsection{Closing routes and factorial cycle moments}
An alternating cycle in the network can be followed by going through one
child and returning through the other. After one child map is fixed, this
produces a permutation containing fresh butterfly layers between fixed
bijections. The following lemma counts several cycles of such a
permutation at once. Its distinguished slots let us sum over all choices
of cycle lengths without losing a factorial.

\begin{lemma}[Factorial moments with distinguished cycle slots]\label{fac:cycle-factorial}
On a finite set take a permutation given by independent butterflies of block size $n=2^r$ in parallel, composed before and after with arbitrary fixed permutations (we allow identifications of the sets at each stage). Fix a subset of size at most $k$, and count cycles lying entirely inside it, with $F_j$ the number of length $j$. For nonnegative integers $a_j$ (finitely supported) and $A=\sum a_j$,
\begin{equation}
\mathbb E\prod_j(F_j)_{a_j}\ \le\ (k/\sqrt n)^A\prod_j j^{-a_j/2}.
\label{fac:e3}
\end{equation}
Here $(F)_a=F(F-1)\cdots(F-a+1)$ and $(F)_0=1$.
\end{lemma}
\begin{proof}
Order the sought cycles by nondecreasing length, with slots distinguished, and choose a starting vertex for each (at most $k^A$ choices). Following a cycle involves querying successive butterfly paths, with the last one required to close the cycle. Demand that these are the specified lengths with no repetitions before closure, are disjoint cycles, and are within the subset; otherwise call this failure. Each closing path has a specified input and output by the time it needs to be queried; conditional on the previous path exposures its probability is zero or $2^u/n$, where $u$ counts switches on that route already exposed (a route through a single butterfly to a given compatible endpoint is unique). If $\mathcal G$ is the exploration just before a closing query, its
required input and output are $\mathcal G$-measurable, and
\[
 \E[\mathbf1_{\{\mathrm{closure}\}}n2^{-u}\mid\mathcal G]\le1,
\]
with value zero for an incompatible route. Ordinary nonclosing queries
insert no weight. Iterating these conditional identities, while
requiring all the distinctness and length conditions, gives
\[
\mathbb E[1_{\rm success} n^A 2^{-U}]\le 1
\]
for each choice of starts, with $U$ the sum of those counts on success. This follows also by sequentially forcing the closing paths when possible, which changes measure by the indicated factors while other path steps cost no weight.

It remains to bound how much the previously exposed paths can help the closing routes. For a realized collection in specified slots, rotate its starts independently uniformly. $U$ counts meetings of selected butterfly paths, namely when the closing path uses a switch met previously. Weight each selected path in a cycle of length $j$ by $1/j$. For meetings between different cycles the expected charge per meeting over rotations is the smaller of the weights (or bounded by it), since cycles are processed by length. Within one cycle it is at most $2/j$. For any subset of $s$ selected paths the total meetings between them in a butterfly layer sequence is at most $sr/2$. Within paths of one cycle of length $j$, Lemma~\ref{net:encounters}
gives the stronger bound $j\log_2j/2$. For several independent
butterfly blocks, apply the lemma in each block and use
$\sum_i j_i\log j_i\le j\log j$ when $\sum_i j_i=j$.
Now integrate the subset bound $sr/2$ over thresholds on the weights: this gives $Ar/2$ for the sum charging each meeting the minimum weight. The additional within-cycle charges give at most $\sum_j a_j\log_2(j)/2$. By Jensen, the sum over rotations ($\prod_j j^{a_j}$ of them) of $2^{-U}$ for each collection is thus at least
\[
\prod_j j^{a_j}\, n^{-A/2}\prod_j j^{-a_j/2}.
\]
To make the counting explicit, let $\mathfrak C(\omega)$ be the set of ordered collections of distinct cycles of the realized permutation, with $a_j$ distinguished slots of length $j$. Its cardinality is $\prod_j(F_j(\omega))_{a_j}$. Each collection has exactly $\prod_jj^{a_j}$ choices of starts. Sum the weighted success inequality over all start choices, then interchange that sum and the expectation. The preceding lower bound for the total rotation weight gives
\[
 k^A\ge \E\sum_{\mathcal C\in\mathfrak C(\omega)}
             \sum_{\text{starts of }\mathcal C} n^A2^{-U}
 \ge n^{A/2}\prod_jj^{a_j/2}\,
              \E\prod_j(F_j)_{a_j}.
\]
Dividing proves \eqref{fac:e3}.
\end{proof}

For later use, if $q\ge1$ and $J\ge1$ is an integer, cycles longer than $J$ contribute at most $k/J$. Expanding $q^{\sum_{j\le J}F_j}$ in factorial moments and using $\sum_{j\le J}j^{-1/2}\le2\sqrt J$ gives
\begin{equation}
\log \mathbb E q^{\sum_j F_j}\le (k/J)\log q+2q k\sqrt J/\sqrt n .
\label{fac:e4}
\end{equation}

\subsection{The two density estimates}\label{net:factorial-route}
\begin{lemma}[Palindromic row-density moments]\label{fac:density}
Let $N=2^d$, $1\le k\le N$, and $a=1/1000$. For a kernel $Q$ on ordered $k$-injections define
\[M(Q,k)=\sup_{x\in\mathcal X_{N,k}}\E_{y\sim u_{N,k}}\big[(N)_kQ(x,y)\big]^{1+a}.\]
There are absolute $C,\eta>0$ with the following properties; below $M$ denotes $M(Q,k)$ for the kernel in question.
\begin{itemize}
\item for $K_N$, $\log M\le C k(k/N)^{\eta}$, for some absolute $C<\infty$ and $\eta>0$;
\item consider the modified palindrome obtained by replacing the central subnetworks of size $S=2^L$ by independent uniform permutations, except in at most $b$ specified such blocks where they are replaced by identity. For every $\xi>0$, there is $L_0$ such that for each fixed $L\ge L_0$ there are $\epsilon>0$ and $d_0\ge L$ for which $M(Q,N)\le e^{\xi N}$ whenever $d\ge d_0$ and $bS/N\le\epsilon$.

\end{itemize}\end{lemma}\begin{proof}
For $N=1$, the only injection has one position and every kernel in the
statement is identity, so $M=1$. Assume henceforth that $d\ge1$.
To see how to estimate these integrals, expand the transition probability by splitting off the outermost layers of each block, recursively. At a given node of size $2^t$, let $h$ be the number of marked paths (the paths of the tuple) there. Given their endpoints at that node, write $b_x,b_y$ for the numbers of occupied direction-$t$ pairs at the two ends. Each term assigns each path to one of the two children, with opposite assignments for a pair coming from the same direction-$t$ edge on either end. Its probability factor for the outer layers is $2^{-b_x-b_y}$. On the labels of paths the two matchings (allowing unmatched vertices) giving these constraints form alternating chains and even cycles, counting parallel matching edges as a cycle as well. If there are $z$ cycles the number of assignments is at most $2^{b_x+b_y-h+z}$, by counting components. Each assignment fixes the child endpoints, and has as further factors the child probabilities. Apply
\[
(\sum_i p_i)^{1+a}\le (\# i)^a\sum_i p_i^{1+a}
\]
at each expansion. On summing over terminal tuple images $y$, keeping one power of each term as the actual routing probability, the remaining factors at expanded nodes give $2^{a(-h+z)}$. Thus, for any initial tuple,
\begin{equation}
\E_{y\sim u_{N,k}} R_x(y)^{1+a}
 \le [(N)_k/N^k]^a\,\mathbb E 2^{a\sum_{t=1}^d Z_t},
\label{fac:e1}
\end{equation}
where $Z_t$ totals the cycles at level $t$ of the actual routing (unused paths can be filled in via the actual switches). For the modified palindrome with $k=N$, stop after levels $d,\ldots,L+1$. Instead the bound is
\begin{equation}
[N!\,(S/N)^N (S!)^{-N/S+b}]^a\,\mathbb E 2^{a\sum_{t>L} Z_t}.
\label{fac:e2}
\end{equation}
Indeed all paths are then used and the extra central probability factors are uniform $(S!)^{-1}$ or identity.

\medskip\noindent\textbf{Fresh layers at one height.}\enspace
Lemma~\ref{fac:cycle-factorial} concerns fresh butterfly layers, whereas
cycle counts at different network heights share coins. Label paths at the
central cross section and read the two butterflies outward, calling them
$X$ and $Y$. Conditioning on $X$ fixes the marked central labels.

Fix a height $t$. Let $\Omega_0,\Omega_1$ be the unions of the bit-zero
and bit-one halves in direction $t$ of all size-$2^t$ nodes. Reading the
first $t-1$ layers of $X$ gives bijections $X_0,X_1$ from these halves to
a common set of switch columns. The input matching, viewed from
$\Omega_0$ to $\Omega_1$, is $I=X_1^{-1}X_0$. The corresponding output
matching is $Y_1^{-1}Y_0$. Hence following an output edge and then an input
edge gives the permutation
\begin{equation}\label{fac:relative-permutation}
 \pi_t=I^{-1}Y_1^{-1}Y_0
       =X_0^{-1}X_1Y_1^{-1}Y_0
       \quad\hbox{on }\Omega_0.
\end{equation}
Each full alternating cycle gives one cycle of $\pi_t$.
The tracked count $Z_t$ includes only those cycles whose vertices on both
sides are marked. In particular it is bounded by the number of cycles of
$\pi_t$ contained in the marked subset of $\Omega_0$, a fixed set of
size at most $k$ under the conditioning on $X$. Switches at layer $t$
change neither matching.

Now expose the first $s<t-1$ outward layers of $Y$, and, for this
single-height estimate, the remaining maps on $\Omega_1$. The map $Y_1$
is then fixed. On $\Omega_0$, write $Y_0=B U$, where $U$ is the exposed
map through height $s$ and $B$ consists of fresh parallel butterflies in
coordinates $s+1,\ldots,t-1$, each of size $n=2^{t-1-s}$. Thus
$\pi_t=(I^{-1}Y_1^{-1})B U$ has exactly the form of
Lemma~\ref{fac:cycle-factorial}. Its bound is uniform in the exposed
maps, so it remains valid after averaging the extra exposure of $Y_1$.

\medskip\noindent\textbf{Summing heights with their conditioning.}\enspace

Here and below group levels above some integer $H\ge1$ into slabs
\[
l<t\le 2l,\qquad l=H,2H,4H,\ldots
\]
cut off at $d$. Put $\mathcal F_s=\sigma(X,\text{the first $s$ outward layers of }Y)$ and
$B_l=\sum_{l<t\le\min(2l,d)}Z_t$. The band variable $B_l$ is
$\mathcal F_{2l}$-measurable (with $\mathcal F_s=\mathcal F_d$ for $s>d$).
Split the slabs into their two alternating families. In either family,
the preceding slab ends at $l/2$, so its cycle count is
$\mathcal F_{\lfloor l/2\rfloor}$-measurable.

Conditioned on that field, the single-level estimate above has a fresh
butterfly of size $n\ge2^{l/2}$. H\"older within the slab uses at most
$l$ factors. For the exponential moment with exponent $2a$, take
$q=2^{2al}$ and $J=\lceil2^{l/16}\rceil$ in \eqref{fac:e4}. Uniformly in
the conditioning,
\begin{equation}
\log\E[q^{Z_t}\mid\mathcal F_{\lfloor l/2\rfloor}]
 \le Ck2^{-c_0l}.
\label{fac:e5}
\end{equation}
The two terms on the right of \eqref{fac:e4} decay exponentially because
$2a<1/16$. Thus, after H\"older,
\[
 \E[2^{2aB_l}\mid\mathcal F_{\lfloor l/2\rfloor}]
 \le\exp(Ck2^{-c_0l}).
\]
For example, if $l_*$ is the last slab in one family and $V$ is the sum
of its earlier bands, the tower property gives
\[
 \E[2^{2a(V+B_{l_*})}\mid X]
 =\E\!\left[2^{2aV}
   \E[2^{2aB_{l_*}}\mid\mathcal F_{\lfloor l_*/2\rfloor}]
   \,\middle|\,X\right]
 \le e^{Ck2^{-c_0l_*}}\E[2^{2aV}\mid X].
\]
Iterate backwards through each family and apply Cauchy--Schwarz to the
two families. The geometric sum of their costs proves
\begin{equation}
\log\E[2^{a\sum_{t>H}Z_t}\mid X]\le Ck2^{-c_1H}.
\label{fac:e6}
\end{equation}
The same reasoning holds with $a$ replaced by $2a$: even the resulting
single-level exponent $4a$ is smaller than $1/16$. We will use this
explicit doubling when separating low and high levels.

\medskip\noindent\textbf{Low heights and the sparse density bound.}\enspace
The high-height cost decays with its starting height. At low heights, sparse occupancy supplies the gain instead. For the low levels of \eqref{fac:e1}, partition the central labels into size $R=2^H$ blocks and let $V_D$ count marked positions in such a block $D$. (Use $H\le d$.) All low-level cycles are internal to such blocks and
\[
\sum_{t\le H} Z_t\le H\sum_D V_D 1_{V_D\ge2}.
\]
For every set of central positions of size $v$ the probability under $X$ they are all marked is at most $p^v$, $p=k/N$. Indeed the upper product-of-marginals bound for every subset holds initially (trace the marked inputs towards the center). It is preserved under a single fair switch: for sets with one endpoint use averaging of the two bounds; with both endpoints use the old product, bounded by the product with both marginals replaced by their mean. After the layers of $X$ each position has marginal $p$.

Write $w=2^{aH}$. In each block use
\[
w^{V_D 1_{V_D\ge 2}}\le 1+\sum_{i=2}^{R}{V_D\choose i} w^i.
\]
Expanding the product and using the joint inclusion bounds yields
\[
\log\mathbb E 2^{a\sum_{t\le H} Z_t}
\le (N/R)\sum_{i=2}^{R}{R\choose i}(p w)^i
\le C(N/R)(R p w)^2
\]
when $Rpw$ is bounded. Choose
$H=\max(1,\lfloor\log_2(1/p)/4\rfloor)$; then $H\le d$.
For the low-height estimate with exponent $2a$, replace $w$ by
$2^{2aH}$. The product $Rp w$ stays bounded, and the resulting logarithmic
cost is at most
\[
 Ck p R2^{4aH}\le C'k p^{(3-4a)/4}.
\]
The high-height estimate \eqref{fac:e6} also holds at exponent $2a$,
with cost at most $Ck2^{-c_1H}$. Low heights still depend on $Y$, so use
Cauchy--Schwarz between the two height ranges, then average over $X$.
Their costs are at most $Ck p^\eta$ for some absolute $\eta>0$.
The normalization in \eqref{fac:e1} is at most one, proving the first
claim. If there are no heights above $H$, only the low-height bound is
needed.

\medskip\noindent\textbf{Uniform central blocks and the dense density bound.}\enspace
For \eqref{fac:e2}, label paths at the input boundary of the central
size-$S$ blocks. Fix the input-side outer butterfly $X$, and include the
uniform or identity central maps at the beginning of $Y$. Use the same
slabs with $H=L$. We distinguish their first members, $l=L$ and $l=2L$,
from the later ones.

At a later slab, $l\ge4L$, condition on $X$, the central maps, and the
ordinary output layers through height $\lfloor l/2\rfloor$. The previous
same-parity slab is measurable in this field. All fresh coordinates used
in \eqref{fac:e5} lie strictly above the central blocks, so that estimate
applies unchanged with $k=N$.

A first slab has no preceding member in its parity family. We therefore
need its moment only conditional on $X$, without fixing the central
permutations. For a height $t$ in this slab expose the bit-one map as in
\eqref{fac:relative-permutation}. In each good bit-zero central block,
realize its uniform permutation as $F U$, where $U$ is an independent
uniform permutation and $F$ is an independent fresh size-$S$ butterfly.
This has the original uniform law. Condition on $U$, but leave $F$ fresh.
In a bad block take $U$ to be identity; the original central map there
is identity. In a comparison experiment insert a fresh $F$ in that block
too, keeping all good-block maps and all later switches the same.

Here is why this insertion costs at most $bS$ cycles. Let $B$ be the
union of the bad block inputs in the fresh-map coordinates. The original
and comparison maps agree on every input outside $B$. After any fixed
premap $U$, the affected domain is $U^{-1}B$, still of size at most $bS$;
a fixed postmap does not enlarge that domain. Every cycle of the
original relative permutation disjoint from $U^{-1}B$ retains all its
arrows and is therefore a cycle of the comparison permutation. At most
$bS$ disjoint cycles meet that domain. Thus, writing $Z'_t$ for the full
comparison cycle count,
\[
 Z_t\le Z'_t+bS.
\]
This comparison holds for every realization of the inserted switches;
no cycle is selected for subsequent conditioning.

In the comparison, the central butterflies and the ordinary output
layers together form fresh parallel butterflies in coordinates
$1,\ldots,t-1$, between fixed maps. Lemma~\ref{fac:cycle-factorial}
therefore bounds their cycle moment. With $q=2^{2al}$, the same choice
$J=\lceil2^{l/16}\rceil$ as before gives, after averaging the auxiliary
maps,
\[
 \log\E[q^{Z_t}\mid X]
 \le CN2^{-c_0l}+bS\log q
 \qquad(l=L\hbox{ or }2L).
\]
H\"older within this first slab costs $O(a l bS)$ in addition to its
decaying term. Integrate the later slabs backwards in each parity family
using their lower-height conditioning, finish with this first-slab
bound, and apply Cauchy--Schwarz between the families. The result is
\[
 \log\mathbb E 2^{a\sum_{t>L}Z_t}
 \le C N 2^{-c_1 L} + C a L bS .
\]
By Stirling, the logarithm of the normalization in square brackets in
\eqref{fac:e2} is at most
$O(\log N+N\log S/S)+b\log(S!)$. Given $\xi>0$, first take $L$
large enough that the terms $N2^{-c_1L}$ and $N\log S/S$ have sufficiently
small coefficients. With that $L$ fixed, choose $\epsilon$ so that
$bS/N\le\epsilon$ controls both bad-block terms, and finally take $d$
large enough to absorb $O(\log N)$. This proves the second claim in the
stated order of choices.
\end{proof}

\subsection{Multiplicity amplification and the middle dimension range}
We now use the branching multiplicity from
\eqref{eq:tuplemultiplicity}: the ordered $k$-injection module contains
$\lambda\vdash N$ with multiplicity $m=f^{\lambda/(N-k)}$, and for
$k=N$ this is $D_\lambda$. The density bounds just proved will make
truncation effective on these repeated copies.

For either symmetric kernel $Q$ above, retain an entry when its density
is at most $z$. Markov's inequality gives discarded row mass at most
$Mz^{-a}$, and symmetry gives the same column bound. The retained
entries are at most $z/(N)_k$. Applying
Lemma~\ref{net:truncation} with $B=z$ therefore gives, on an irreducible
occurring with multiplicity $m>0$,
\begin{equation}
 \|Q(\lambda)\|\le Mz^{-a}+\sqrt{z/m}.
 \label{fac:e7}
\end{equation}
In particular if $\log M\le (a/4)\log m$, take $z=\sqrt m$ to get at most $2m^{-a/4}$.

First consider $j=N-\lambda_1$ with $1\le j\le\delta N$, where a sufficiently small absolute $\delta>0$ will be used. Choose a fixed $u_0>0$ small enough that $\eta(1-u_0)>u_0$ and $u_0<1$, and take $k=\lceil j(N/j)^{u_0}\rceil$. For small $\delta$ this gives $N-k\ge j$, so the skew tableau of shape $\lambda/(N-k)$ has its remaining first row independent of the entire tail shape $\rho=(\lambda_2,\ldots)$. Consequently
\[
m=\binom{k}{j} f^\rho,\qquad D_\lambda\le \binom Nj f^\rho.
\]
Thus $\log m\ge u_0 j\log(N/j)+\log f^\rho$, at least a fixed positive fraction of $\log D_\lambda$ (using the elementary binomial bounds). The first density bound proved above gives $\log M\le C' j(j/N)^{\eta(1-u_0)-u_0}$, hence at most $(a/4)\log m$ by reducing $\delta$. By \eqref{fac:e7}, and $m$ tending uniformly to infinity here as $d\to\infty$, we have $\|K_N(\lambda)\|\le D_\lambda^{-c_2}$ for some constant $c_2>0$.

By Lemma~\ref{lem:annihilation}, shapes with more than $N/2$ rows
are annihilated by $T_N$ and hence by $K_N$.

All remaining shapes other than the trivial one and those already treated have $j\ge\delta N$ and
\[
\log D_\lambda\ge hN
\]
with $h>0$ fixed, for large $N$. Here are details. Both first row and column are then bounded by $(1-\delta)N$ (take $\delta<1/2$). If both are smaller than $N/10$, hook lengths are at most $N/5$ and the bound follows from Stirling. Otherwise transpose if necessary (which does not change dimension) so the row has length $s\ge N/10$. Retain only it and a subdiagram of the tail of size $v=\lfloor\min(\delta/2,1/100)N\rfloor$, which can be obtained by removing corners. All tableaux on this diagram have extensions. Fill its first $v$ positions in the top row first, then interleave the remaining row and a fixed tableau order on the tail; this gives $\binom{s}{v}$ possibilities, yielding the asserted exponential lower bound.

Take a fixed large $A_0$ so that when $\log D_\lambda\ge A_0 N$, the first density bound on all $N$ positions gives $\log M\le(a/4)\log D_\lambda$. Then \eqref{fac:e7} with $Q=K_N$ again gives $\|K_N(\lambda)\|\le D_\lambda^{-c_3}$ for an absolute $c_3>0$, for large $N$.

The only remaining range is $hN\le\log D_\lambda\le A_0 N$. Here the full-deck moment cost $CN$ is comparable with the log dimension and need not fit the truncation threshold. We use the modified central networks to reduce that cost before invoking multiplicity. Choose $L$ fixed sufficiently large and then $\epsilon>0$ sufficiently small so the second of the two assertions about the $(1+a)$-moment above gives $\log M\le ahN/4$ whenever $b\le\epsilon N/S$, for large $N$. We justify the comparison in positive semidefinite order. Let $O$ be
the product of the outer butterfly layers, and let $C$ be the parallel
product of the central size-$S$ palindrome operators. Thus, on the
representation under consideration,
\[
 K_N=O^*CO.
\]
By Lemma~\ref{lem:finitegap}, each central positive contraction has norm
at most a fixed $\theta<1$ on every nontrivial irreducible of $S_S$;
increase $\theta$ to a number in $(0,1)$ if necessary.

Restrict to $(S_S)^{N/S}$ and decompose into tensor-product irreducibles,
with their multiplicity spaces. If a component is nontrivial on more
than $\epsilon N/S$ blocks, the tensor product defining $C$ has norm at
most $\theta^{\epsilon N/S}$. This accounts for the first term below.

For $E\subseteq[N/S]$, let $P_E$ act as identity in blocks in $E$ and
as the projection to invariants in the other blocks. On a component
whose nontrivial factors lie in a set $F$ with $|F|\le\epsilon N/S$,
$P_F$ is identity. Consequently $\sum_{|E|\le\epsilon N/S}P_E$
dominates identity on all these components and therefore dominates
$C$ there. Its sandwich
\[
 Q_E=O^*P_EO
\]
is precisely the modified network with identity in $E$ and uniform
central permutations elsewhere. Since $O$ is a contraction, these two
component cases give
\[
0\le K_N(\lambda)\le \theta^{\epsilon N/S} I+\sum_{E:\, |E|\le\epsilon N/S}Q_E(\lambda)
\]
where $Q_E$ is exactly the modified palindrome with identity in $E$ and uniform central permutations elsewhere. Each summand by \eqref{fac:e7}, $k=N$, is at most $2D_\lambda^{-a/4}$ in norm. The log number of summands is at most $(N/S)[-\epsilon\log\epsilon-(1-\epsilon)\log(1-\epsilon)]$ by the binomial theorem ($\epsilon<1/2$), so we may require also that this be $\le ah N/8$. Since $hN\le\log D_\lambda\le A_0 N$ and all parameters now are fixed absolutely, this proves $\|K_N(\lambda)\|\le D_\lambda^{-c_4}$ for some absolute $c_4>0$ and sufficiently large $N$.

We have consequently, on every nontrivial shape not killed, $\|T_N(\lambda)\|\le D_\lambda^{-c_*}$ for an absolute $c_*>0$. Plancherel on the finite group and Cauchy-Schwarz bound four times the squared total variation distance after $r$ butterflies by
\[
\sum_{\lambda\ne (N)} D_\lambda\,\|T_N(\lambda)^r\|_{\rm HS}^2,
\]
independently of initial permutation by translation. For $r$ forward sweeps,
\[
 \|T_N(\lambda)^r\|_{\HS}^2
 \le D_\lambda\|T_N(\lambda)^r\|_{\op}^2
 \le D_\lambda\|T_N(\lambda)\|_{\op}^{2r}.
\]
Thus, keeping only shapes not killed, the sum is at most
$\sum D_\lambda^{2-2rc_*}$.

This last bound tends to zero for sufficiently large fixed $r$. Medium or larger dimensions with $\log D_\lambda\ge hN$ have at most $2^N$ terms (bound partitions by compositions). For the small-$j$ range, $D_\lambda\ge \binom{N-j}{j}$: again fill $j$ positions of the first row first and then interleave tail and remaining row. There are at most $2^j$ shapes per $j$. Thus, with $s_*=2r c_*-2>0$, this range contributes at most
\[
\sum_{1\le j\le\delta N}2^j(j/(N-j))^{s_*j}
\]
tending to zero (choose $s_*$ large for a geometric summable majorant and each fixed term vanishes). Increase $r$ also so that $s_* h>\log 2$. Hence $rd$ steps suffice for the required threshold for all sufficiently large $d$.

